\documentclass[11pt,a4paper]{article}
\usepackage[T1]{fontenc}
\usepackage[margin=2.5cm]{geometry}
\usepackage{amsmath,amssymb,amsthm}
\usepackage{enumitem}
\usepackage[colorlinks=true,linkcolor=blue,urlcolor=blue,citecolor=blue]{hyperref}
\setlength{\emergencystretch}{2em}

\newtheorem{theorem}{Theorem}[section]
\newtheorem{proposition}[theorem]{Proposition}
\newtheorem{lemma}[theorem]{Lemma}
\newtheorem{corollary}[theorem]{Corollary}
\theoremstyle{remark}
\newtheorem{remark}[theorem]{Remark}

\newcommand{\ZZ}{\mathbb{Z}}
\newcommand{\FF}{\mathbb{F}}
\newcommand{\one}{\mathbf{1}}
\newcommand{\sg}{\sigma}
\newcommand{\Nb}{\Gamma}

\title{Odd eigenfunctions of squared norm $16$ on a strongly regular graph $(99,14,1,2)$ with an involution\\[4pt]
\large Part II.A of the exclusion of automorphisms of order two: the computer-free steps of the first lemma}
\author{Working draft, 14 October 2026}
\date{}

\begin{document}
\maketitle

\begin{abstract}
Let $\Gamma$ be a strongly regular graph with parameters $(99,14,1,2)$, $\sg$ an automorphism of order $2$ and $A$ the
adjacency matrix. The first of the three lemmas used in the second half of the exclusion of $\sg$ says that there is
no integer-valued function $\varphi$ with $\varphi\circ\sg=-\varphi$, $A\varphi=-4\varphi$ and
$\sum_x\varphi(x)^2=16$. This text contains, with every verification, the part of its proof that needs no computer:
such a function takes the values $0,\pm1$, its support avoids the neighbours of the fixed vertex, and the bipartite
graph between its positive and negative supports is the complement of the cube or the complement of the Wagner
graph; the cube is impossible. The complement of the Wagner graph is the one case that the existing proof settles by
a (small) exhaustive enumeration; it is not treated here.
\end{abstract}

%======================================================================
\section{Setting and statement}

We use the notation and the results of Part~I: $\Gamma$ is a strongly regular graph with parameters $(99,14,1,2)$;
$x\sim y$ means that $x,y$ are adjacent and $\Nb(x)$ is the set of neighbours of $x$; two adjacent vertices have
exactly one common neighbour and two distinct non-adjacent vertices exactly two; $\sg$ is an automorphism of order
$2$. We use three facts proved in Part~I.
\begin{enumerate}[label=\textup{(F\arabic*)},leftmargin=3em]
\item\label{F1} $\sg$ fixes exactly one vertex $x_0$. The $14$ neighbours of $x_0$ (the \emph{inner} vertices) form
seven pairs $\{a,\sg a\}$ of adjacent vertices, the \emph{blocks}; two inner vertices in different blocks are not
adjacent. The $84$ other vertices different from $x_0$ are called \emph{outer}.
\item\label{F2} A vertex $v$ with $v\sim\sg v$ is inner.
\item\label{F3} An outer vertex has exactly two inner neighbours, in two different blocks.
\end{enumerate}
A function $\varphi$ on the vertices is \emph{odd} if $\varphi\circ\sg=-\varphi$; then $\varphi(x_0)=0$.

\begin{theorem}\label{thm:A}
Let $\varphi$ be an odd function with integer values such that $A\varphi=-4\varphi$ and $\sum_x\varphi(x)^2=16$. Put
$P=\{x:\varphi(x)=1\}$ and $Q=\{x:\varphi(x)=-1\}$. Then:
\begin{enumerate}[label=\textup{(\alph*)}]
\item $\varphi$ takes only the values $0,1,-1$; $|P|=|Q|=8$ and $Q=\sg P$;
\item no two vertices of $P$ are adjacent, no two vertices of $Q$ are adjacent, and every vertex of $P$ has exactly
four neighbours in $Q$;
\item all vertices of $P\cup Q$ are outer;
\item the graph $\mathcal C$ on $P$ in which two distinct vertices $x,y$ are joined when $x\not\sim\sg y$ is
isomorphic to the Wagner graph (the cycle of length $8$ with its four long diagonals).
\end{enumerate}
\end{theorem}

The graph $\mathcal C$ is well defined as an undirected graph because $x\sim\sg y$ if and only if $\sg x\sim y$.
Theorem~\ref{thm:A} leaves one configuration; that it does not occur either is the remaining part of the lemma.

%======================================================================
\section{Values, independence, and the matrix $B$}

Throughout, $\varphi$ is as in Theorem~\ref{thm:A}.

\begin{lemma}\label{lem:values}
$\varphi$ takes only the values $0,1,-1$.
\end{lemma}
\begin{proof}
A value of absolute value at least $3$ at $x$ would give $\varphi(x)^2+\varphi(\sg x)^2\ge18>16$. Suppose
$|\varphi(x)|=2$. Then $\varphi(\sg x)=-\varphi(x)$, and the other vertices satisfy
$\sum_{y\ne x,\sg x}\varphi(y)^2=8$. The equation $(A\varphi)(x)=-4\varphi(x)$ reads
$\sum_{y\sim x}\varphi(y)=-4\varphi(x)$, of absolute value $8$. If $x\sim\sg x$, the neighbour $\sg x$ contributes
$-\varphi(x)$ and the other neighbours must contribute $-3\varphi(x)$, of absolute value $6$; otherwise they
contribute $-4\varphi(x)$, of absolute value $8$. Now the vertices $y\ne x,\sg x$ come in pairs $\{y,\sg y\}$ with
$\varphi(\sg y)=-\varphi(y)$ (and $\varphi(x_0)=0$). A pair with both vertices adjacent to $x$ contributes $0$ to
the sum; a pair with one vertex adjacent to $x$ contributes $\pm\varphi(y)$. So the absolute value of the
contribution of the neighbours other than $\sg x$ is at most $\sum|\varphi(y)|$ over one vertex $y$ per pair, and
$\sum|\varphi(y)|\le\sum\varphi(y)^2=4$ over these representatives. This is less than $6$. Contradiction.
\end{proof}

By Lemma~\ref{lem:values}, $\sum_x\varphi(x)^2=16$ says $|P|+|Q|=16$, and $Q=\sg P$ since $\varphi$ is odd. This is
Theorem~\ref{thm:A}(a). Since $\varphi$ vanishes outside $P\cup Q$, for every vertex $z$,
\begin{equation}\label{eq:eig}
(A\varphi)(z)=|\Nb(z)\cap P|-|\Nb(z)\cap Q| .
\end{equation}

\begin{lemma}\label{lem:indep}
No two vertices of $P$ are adjacent, and every vertex of $P$ has exactly four neighbours in $Q$. The same holds with
$P$ and $Q$ exchanged.
\end{lemma}
\begin{proof}
For $x\in P$ let $d_x=|\Nb(x)\cap P|$, and let $e=\frac12\sum_{x\in P}d_x$ be the number of edges inside $P$. By
\eqref{eq:eig} and $(A\varphi)(x)=-4$, $x$ has $d_x+4$ neighbours in $Q$.

Count the triples $(\{x,y\},z)$ with $x\ne y$ in $P$ and $z$ a common neighbour of $x$ and $y$. Adjacent pairs have
one common neighbour and the others two, so there are $2\binom82-e=56-e$ triples. A vertex $z$ is in
$\binom{n_z}2$ triples, where $n_z=|\Nb(z)\cap P|$. For $z\in P$, $n_z=d_z$. For $z=\sg x\in Q$, the neighbours of
$\sg x$ in $P$ are the images under $\sg$ of the neighbours of $x$ in $Q$, so $n_z=d_x+4$. The vertices of $P\cup Q$
are therefore in
\[
\sum_{x\in P}\Bigl[\binom{d_x}2+\binom{d_x+4}2\Bigr]=\sum_{x\in P}(d_x^2+3d_x+6)=\sum_{x\in P}d_x^2+6e+48
\]
triples, and the number of triples with $z\notin P\cup Q$ is
\begin{equation}\label{eq:D}
D=56-e-\sum_{x\in P}d_x^2-6e-48=8-7e-\sum_{x\in P}d_x^2\ \ge0 .
\end{equation}
Since $d_x^2\ge d_x$ and $\sum_xd_x=2e$, we get $0\le D\le8-9e$, so $e=0$ and all $d_x=0$. The statement for $Q$
follows by applying $\sg$.
\end{proof}

This is Theorem~\ref{thm:A}(b). We record the value $D=8$ in \eqref{eq:D}:
\begin{equation}\label{eq:D8}
\text{exactly $8$ triples $(\{x,y\},z)$ with $x\ne y$ in $P$ have their common neighbour $z$ outside $P\cup Q$.}
\end{equation}

\paragraph{The matrix $B$.} Let $B$ be the matrix with rows and columns indexed by $P$ and
\[
B[x,y]=\begin{cases}1&\text{if }x\sim\sg y,\\0&\text{otherwise.}\end{cases}
\]
For $x\in P$ we call $\mathrm{row}(x)=\{y\in P:B[x,y]=1\}$ the \emph{row} of $x$.

\begin{lemma}\label{lem:B}
\begin{enumerate}[label=\textup{(\alph*)}]
\item $B$ is symmetric, and every row has exactly four elements.
\item $x\in\mathrm{row}(x)$ if and only if $x$ is inner.
\item For $x\ne y$ in $P$, $|\mathrm{row}(x)\cap\mathrm{row}(y)|$ is the number of common neighbours of $x$ and $y$
in $Q$; it is at most $2$, and at most $1$ if $x$ and $y$ have a common neighbour outside $Q$.
\end{enumerate}
\end{lemma}
\begin{proof}
(a) $x\sim\sg y$ if and only if $\sg x\sim y$, that is, $y\sim\sg x$. The row of $x$ is in bijection with
$\Nb(x)\cap Q$ by $y\mapsto\sg y$, and $|\Nb(x)\cap Q|=4$ by Lemma~\ref{lem:indep}.
(b) $B[x,x]=1$ means $x\sim\sg x$; by \ref{F2} this implies that $x$ is inner, and conversely an inner vertex is
adjacent to its image by \ref{F1}.
(c) $y'\in\mathrm{row}(x)\cap\mathrm{row}(y)$ means that $\sg y'\in Q$ is adjacent to $x$ and to $y$. By
Lemma~\ref{lem:indep}, $x\not\sim y$, so $x$ and $y$ have exactly two common neighbours.
\end{proof}

%======================================================================
\section{The support avoids the neighbours of the fixed vertex}

Let $k$ be the number of inner vertices in $P$. Since each block is an orbit $\{a,\sg a\}$, a block contains a
vertex of $P$ if and only if it contains a vertex of $Q$, and then exactly one of each; so $k$ is the number of
blocks meeting $P\cup Q$, and two inner vertices of $P$ lie in different blocks.

\begin{lemma}\label{lem:k02}
$k\in\{0,2\}$.
\end{lemma}
\begin{proof}
\emph{$k$ is even.} The sum of all entries of $B$ is $8\cdot4=32$ by Lemma~\ref{lem:B}(a). Since $B$ is symmetric
with entries $0$ and $1$, this sum is the number of $x$ with $B[x,x]=1$ plus twice the number of unordered pairs
$\{x,y\}$, $x\ne y$, with $B[x,y]=1$. By Lemma~\ref{lem:B}(b) the first number is $k$. So $k$ is even.

\emph{$k\le3$.} Let $a\in P$ be inner. Its $14$ neighbours are $x_0$, the vertex $\sg a\in Q$ and $12$ outer
vertices. By Lemma~\ref{lem:indep}, $a$ has exactly four neighbours in $Q$; so exactly three outer vertices of $Q$
are adjacent to $a$. Count the pairs $(a,y)$ with $a\in P$ inner, $y\in Q$ outer and $a\sim y$: there are $3k$ of
them. On the other hand there are $8-k$ outer vertices in $Q$, and each has exactly two inner neighbours by
\ref{F3}. Hence $3k\le2(8-k)$, that is, $5k\le16$.
\end{proof}

\begin{lemma}\label{lem:k0}
$k=0$.
\end{lemma}
\begin{proof}
By Lemma~\ref{lem:k02} it suffices to exclude $k=2$. Suppose $k=2$ and let $p_0,p_1$ be the two inner vertices of
$P$; they lie in different blocks. By
Lemma~\ref{lem:B}(b), $p_0\in\mathrm{row}(p_0)$, $p_1\in\mathrm{row}(p_1)$, and $x\notin\mathrm{row}(x)$ for the six
other vertices $x$ of $P$; let $R$ be the set of these six.

\emph{$p_1\notin\mathrm{row}(p_0)$}: $p_0$ and $\sg p_1$ are inner vertices in different blocks, so they are not
adjacent by \ref{F1}.

\emph{$|\mathrm{row}(p_0)\cap\mathrm{row}(p_1)|\le1$}: $p_0$ and $p_1$ have the common neighbour $x_0$, which is not
in $Q$; apply Lemma~\ref{lem:B}(c).

So $\mathrm{row}(p_0)=\{p_0\}\cup A_0$ and $\mathrm{row}(p_1)=\{p_1\}\cup C_0$ with $A_0,C_0\subset R$,
$|A_0|=|C_0|=3$ and $|A_0\cap C_0|\le1$.

\medskip\noindent\emph{Case $A_0\cap C_0=\emptyset$.} Let $a\in A_0$. By symmetry of $B$, $p_0\in\mathrm{row}(a)$
and $p_1\notin\mathrm{row}(a)$; also $a\notin\mathrm{row}(a)$. So $\mathrm{row}(a)$ consists of $p_0$ and three
elements of $(A_0\setminus\{a\})\cup C_0$. Now $\mathrm{row}(a)\cap\mathrm{row}(p_0)$ contains $p_0$ and
$\mathrm{row}(a)\cap A_0$; by Lemma~\ref{lem:B}(c), $|\mathrm{row}(a)\cap A_0|\le1$. And
$\mathrm{row}(a)\cap\mathrm{row}(p_1)=\mathrm{row}(a)\cap C_0$ has at most two elements. Hence
$|\mathrm{row}(a)\cap A_0|=1$ for every $a\in A_0$: the relation ``$a'\in\mathrm{row}(a)$'' is a graph on the three
elements of $A_0$ in which every vertex has degree $1$. This is impossible, since the sum of the degrees is even.

\medskip\noindent\emph{Case $|A_0\cap C_0|=1$.} Write $A_0=\{s,a,b\}$, $C_0=\{s,c,d\}$ and let $t$ be the sixth
element of $R$.

(1) \emph{$\mathrm{row}(t)=\{a,b,c,d\}$.} Since $t\notin A_0\cup C_0$, $p_0,p_1\notin\mathrm{row}(t)$, and
$t\notin\mathrm{row}(t)$; so $\mathrm{row}(t)$ consists of four of the five elements $s,a,b,c,d$. By
Lemma~\ref{lem:B}(c), $|\mathrm{row}(t)\cap A_0|\le2$ and $|\mathrm{row}(t)\cap C_0|\le2$. If $s\in\mathrm{row}(t)$,
the sum of these two numbers would be $2+3=5$. So $s\notin\mathrm{row}(t)$.

(2) \emph{After exchanging $a,b$ and $c,d$ if necessary, $\mathrm{row}(s)=\{p_0,p_1,a,c\}$.} We have
$p_0,p_1\in\mathrm{row}(s)$, $s\notin\mathrm{row}(s)$, and $t\notin\mathrm{row}(s)$ by (1) and symmetry. So
$\mathrm{row}(s)$ contains two elements of $\{a,b,c,d\}$. As $\mathrm{row}(s)\cap\mathrm{row}(p_0)$ contains $p_0$
and $\mathrm{row}(s)\cap\{a,b\}$, at most one of $a,b$ is in $\mathrm{row}(s)$; similarly at most one of $c,d$.

(3) \emph{$b\notin\mathrm{row}(a)$ and $d\notin\mathrm{row}(c)$.} $\mathrm{row}(a)$ contains $p_0$ and $s$ (by (2)
and symmetry), which both belong to $\mathrm{row}(p_0)=\{p_0,s,a,b\}$; by Lemma~\ref{lem:B}(c), $b\notin\mathrm{row}(a)$.
The same argument with $p_1$ gives $d\notin\mathrm{row}(c)$.

(4) \emph{$\mathrm{row}(b)=\{p_0,t,c,d\}$ and $\mathrm{row}(d)=\{p_1,t,a,b\}$.} $\mathrm{row}(b)$ contains $p_0$ and
$t$ (by (1)); it does not contain $b$, $p_1$ (as $b\notin C_0$), $s$ (by (2)) or $a$ (by (3)). Only $c,d$ remain,
and the row has four elements. Similarly $\mathrm{row}(d)$ contains $p_1$ and $t$, and does not contain $d$, $p_0$,
$s$ or $c$; so it contains $a$ and $b$.

(5) \emph{Contradiction.} $\mathrm{row}(a)$ contains $p_0$, $s$, $t$ (by (1)) and $d$ (by (4)); so
$\mathrm{row}(a)=\{p_0,s,t,d\}$. Then $\mathrm{row}(a)\cap\mathrm{row}(b)\supseteq\{p_0,t,d\}$ has three elements,
contradicting Lemma~\ref{lem:B}(c).
\end{proof}

This is Theorem~\ref{thm:A}(c): all vertices of $P\cup Q$ are outer, and $B$ has zero diagonal.

%======================================================================
\section{The complement of $B$ is the cube or the Wagner graph}

By Lemma~\ref{lem:k0} and Lemma~\ref{lem:B}, the graph $\mathcal C$ of Theorem~\ref{thm:A}(d) has adjacency matrix
$J-I-B$, and every vertex of $\mathcal C$ has $7-4=3$ neighbours.

\begin{lemma}\label{lem:C}
$\mathcal C$ has no triangle. Moreover, for $x\ne y$ in $P$, the number of common neighbours of $x$ and $y$ in $Q$
is $2$ if $x,y$ are adjacent in $\mathcal C$, and equals the number of common neighbours of $x,y$ in $\mathcal C$
otherwise.
\end{lemma}
\begin{proof}
Let $M=J-I-B$. Then $B=J-I-M$, $J^2=8J$ and $JM=MJ=3J$, so
\[
B^2=J^2-2J+I-JM-MJ+2M+M^2=I+2M+M^2 .
\]
For $x\ne y$ the entry $(x,y)$ of $B^2$ is $|\mathrm{row}(x)\cap\mathrm{row}(y)|$, the number of common neighbours
of $x,y$ in $Q$ (Lemma~\ref{lem:B}(c)), and it is at most $2$. So $2M[x,y]+M^2[x,y]\le2$: if $x,y$ are adjacent in
$\mathcal C$ they have no common neighbour in $\mathcal C$ and the count is $2$; otherwise the count is $M^2[x,y]$.
\end{proof}

\begin{lemma}\label{lem:cubic}
A graph on $8$ vertices in which every vertex has degree $3$ and which has no triangle is isomorphic to the cube
$Q_3$ or to the Wagner graph $M_8$ (vertices $0,\dots,7$ modulo $8$, $i$ adjacent to $i\pm1$ and to $i+4$).
\end{lemma}
\begin{proof}
Let $v$ be a vertex, $a,b,c$ its neighbours and $R$ the set of the four other vertices. As there is no triangle,
$a,b,c$ are pairwise non-adjacent, so each has two neighbours in $R$: six edges between $\{a,b,c\}$ and $R$. The
sum of the degrees of the vertices of $R$ is $12$, so $R$ contains $(12-6)/2=3$ edges. A graph with four vertices,
three edges and no triangle is a star (one vertex adjacent to the three others) or a path with four vertices.

\emph{Star.} Let $r_0$ be the centre. It has no neighbour among $a,b,c$, and each of the three leaves has two.
Each of $a,b,c$ has two neighbours among the leaves. The bipartite graph between $\{a,b,c\}$ and the leaves has all
degrees $2$ and six vertices, so it is a cycle of length $6$. The graph is then determined up to isomorphism; it is
the cube, with $v$ and $r_0$ opposite.

\emph{Path.} Let $R$ be the path $r_0,r_1,r_2,r_3$. Then $r_0,r_3$ have two neighbours among $a,b,c$ and $r_1,r_2$
have one. The two neighbours in $R$ of each of $a,b,c$ are not adjacent, so they form one of the pairs
$\{r_0,r_2\}$, $\{r_0,r_3\}$, $\{r_1,r_3\}$. If these pairs are used $n_{02},n_{03},n_{13}$ times, the degrees give
$n_{13}=1$ (for $r_1$), $n_{02}=1$ (for $r_2$), and $n_{02}+n_{03}=2$ (for $r_0$), so $n_{03}=1$. Up to the names of
$a,b,c$: $a\sim r_0,r_2$; $b\sim r_0,r_3$; $c\sim r_1,r_3$. The map
\[
v\mapsto0,\quad c\mapsto1,\quad r_1\mapsto2,\quad r_0\mapsto3,\quad b\mapsto4,\quad r_3\mapsto5,\quad
r_2\mapsto6,\quad a\mapsto7
\]
is an isomorphism onto $M_8$: the edges $vc$, $cr_1$, $r_1r_0$, $r_0b$, $br_3$, $r_3r_2$, $r_2a$, $av$ become the
cycle $0,1,\dots,7$, and the edges $vb$, $cr_3$, $r_1r_2$, $r_0a$ become $\{0,4\}$, $\{1,5\}$, $\{2,6\}$, $\{3,7\}$.
\end{proof}

%======================================================================
\section{The vertices outside the support}

For a vertex $z\notin P\cup Q$, \eqref{eq:eig} and $(A\varphi)(z)=0$ give $|\Nb(z)\cap P|=|\Nb(z)\cap Q|$; we call
this number $d_z$, and we let $X_d$ be the set of $z\notin P\cup Q$ with $d_z=d$. Since $\sg$ exchanges $P$ and $Q$,
$d_{\sg z}=d_z$: each $X_d$ is mapped to itself by $\sg$. The fixed vertex $x_0$ is in $X_0$, because its neighbours
are inner and $P\cup Q$ consists of outer vertices.

\begin{lemma}\label{lem:X}
$d_z\le2$ for every $z\notin P\cup Q$, and $|X_0|=11$, $|X_1|=64$, $|X_2|=8$.
\end{lemma}
\begin{proof}
Suppose $z$ is adjacent to three vertices $x_1,x_2,x_3$ of $P$. They are pairwise non-adjacent, so two of them have
exactly two common neighbours, one of which is $z\notin Q$; hence their sets of neighbours in $Q$, of size $4$ each,
meet pairwise in at most one vertex. The union of these three sets has at least $12-3=9$ elements, but $|Q|=8$.

Each $x\in P$ has $14-4=10$ neighbours outside $P\cup Q$, so $\sum_{z\notin P\cup Q}d_z=80$. By \eqref{eq:D8},
$\sum_{z\notin P\cup Q}\binom{d_z}2=8$; as $d_z\le2$ this says $|X_2|=8$. Then $|X_1|=80-2\cdot8=64$ and
$|X_0|=99-16-8-64=11$.
\end{proof}

For $z\notin P\cup Q$ let $\alpha_z$, $\beta_z$, $\gamma_z$ be the numbers of neighbours of $z$ in $X_2$, $X_1$,
$X_0$.

\begin{lemma}\label{lem:abc}
For $z\in X_d$: $\ \beta_z+2\alpha_z=16-5d$ and $\gamma_z=\alpha_z+3d-2$. In particular $\alpha_z\ge2$ and
$\gamma_z=\alpha_z-2$ for $z\in X_0$, and $\gamma_z=\alpha_z+4$ for $z\in X_2$.
\end{lemma}
\begin{proof}
Count the pairs $(w,s)$ with $w\sim z$, $s\in P\cup Q$ and $w\sim s$. For a fixed $s$, the number of $w$ is the
number of common neighbours of $z$ and $s$: $1$ for the $2d$ vertices $s$ adjacent to $z$ and $2$ for the
$16-2d$ others; total $2d+2(16-2d)=32-2d$. For a fixed $w\in\Nb(z)$, the number of $s$ is the number of neighbours of
$w$ in $P\cup Q$: $4$ if $w\in P\cup Q$ (Lemma~\ref{lem:indep}), $2d'$ if $w\in X_{d'}$. So
$4\cdot2d+4\alpha_z+2\beta_z=32-2d$, which is the first identity. The degree of $z$ is
$14=2d+\alpha_z+\beta_z+\gamma_z$; substituting $\beta_z=16-5d-2\alpha_z$ gives the second.
\end{proof}

\begin{corollary}\label{cor:ten}
$\sum_{z\in X_0}(\alpha_z-2)\ge10$.
\end{corollary}
\begin{proof}
Let $\tau$ be the number of edges inside $X_2$, so $\sum_{z\in X_2}\alpha_z=2\tau$. The number of edges between
$X_0$ and $X_2$ is $\sum_{z\in X_2}\gamma_z=\sum_{z\in X_2}(\alpha_z+4)=2\tau+32$, and also
$\sum_{z\in X_0}\alpha_z$. As $|X_0|=11$, $\sum_{z\in X_0}(\alpha_z-2)=2\tau+32-22=10+2\tau$.
\end{proof}

%======================================================================
\section{The cube is impossible}

\begin{proposition}\label{prop:cube}
$\mathcal C$ is not isomorphic to the cube.
\end{proposition}
\begin{proof}
Suppose it is, and label the vertices of $P$ by the elements of $\FF_2^3$ so that two vertices are adjacent in
$\mathcal C$ exactly when their labels differ in one coordinate. Two distinct vertices of the cube at distance $1$
are adjacent; at distance $2$ they have two common neighbours; at distance $3$ (\emph{antipodal} vertices) they have
none. By Lemma~\ref{lem:C}, two distinct vertices of $P$ have two common neighbours in $Q$ unless they are antipodal,
in which case they have none. Recall that two distinct vertices of $P$ have exactly two common neighbours in
$\Gamma$, none of them in $P$ (Lemma~\ref{lem:indep}). Hence:
\begin{equation}\label{eq:anti}
\begin{minipage}{0.86\textwidth}
a vertex outside $P\cup Q$ adjacent to two distinct vertices $x,y$ of $P$ exists only if $x,y$ are antipodal, and
for each of the four antipodal pairs there are exactly two such vertices.
\end{minipage}
\end{equation}

\medskip\noindent\emph{Step 1: the vertices of $X_2$ are the arcs of a $4$-cycle.}
Let $T_0,T_1,T_2,T_3$ be the four antipodal pairs of $P$. By \eqref{eq:anti}, for every $z\in X_2$ the set
$\Nb(z)\cap P$ is one of the $T_i$, and each $T_i$ arises from exactly two vertices of $X_2$. For $z\in X_2$ let
$i(z)$ be the index with $\Nb(z)\cap P=T_{i(z)}$, and $j(z)=i(\sg z)$; then
$\Nb(z)\cap Q=\sg\bigl(\Nb(\sg z)\cap P\bigr)=\sg T_{j(z)}$. Note that $z\ne\sg z$ for $z\in X_2$, since
$x_0\in X_0$; and $i(\sg z)=j(z)$, $j(\sg z)=i(z)$.

\emph{$i(z)\ne j(z)$.} Otherwise $z$ and $\sg z$ would both be adjacent to the two vertices of $T_{i(z)}$ and to the
two vertices of $\sg T_{i(z)}$: two distinct vertices with four common neighbours.

\emph{The map $z\mapsto(i(z),j(z))$ is injective on $X_2$.} Two distinct vertices of $X_2$ with the same image
would have the same four neighbours in $P\cup Q$.

So $X_2$ is identified with a set of $8$ ordered pairs $(i,j)$, $i\ne j$, closed under $(i,j)\mapsto(j,i)$, in
which each $i$ occurs exactly twice as first entry. The $4$ unordered pairs $\{i,j\}$ form a graph without loops
or repeated edges on $\{0,1,2,3\}$ in which every vertex has degree $2$: a cycle of length $4$. Renumbering the
$T_i$, we may assume that it is the cycle $0,1,2,3$; so
\[
X_2=\{z_{ij}:\{i,j\}\in\{\{0,1\},\{1,2\},\{2,3\},\{3,0\}\}\},\qquad
\Nb(z_{ij})\cap P=T_i,\quad\Nb(z_{ij})\cap Q=\sg T_j .
\]

\medskip\noindent\emph{Step 2: two cycles of length $4$ on $X_2$.}
Two distinct vertices $z_{ij}$, $z_{i'j'}$ have, as common neighbours in $P\cup Q$, the two vertices of $T_i$ if
$i=i'$, the two vertices of $\sg T_j$ if $j=j'$, and none otherwise (the $T_i$ are pairwise disjoint; $i=i'$ and
$j=j'$ do not both hold). Call them \emph{linked} when $i=i'$ or $j=j'$. The four vertices
$z_{01},z_{03},z_{21},z_{23}$ are linked along the cycle
\[
z_{01}\ -\ z_{03}\ -\ z_{23}\ -\ z_{21}\ -\ z_{01}
\]
(common first index, common second index, alternately), and $z_{01},z_{23}$ are not linked, nor are
$z_{03},z_{21}$. The four vertices $z_{10},z_{30},z_{12},z_{32}$ are linked along the cycle
$z_{10}-z_{30}-z_{32}-z_{12}-z_{10}$ in the same way. A vertex of the first group and a vertex of the second are not
linked: their first indices have different parities, and so have their second indices. So ``linked'' consists of two
disjoint cycles of length $4$. We call \emph{diagonals} the four pairs
$\{z_{01},z_{23}\}$, $\{z_{03},z_{21}\}$, $\{z_{10},z_{32}\}$, $\{z_{30},z_{12}\}$ of opposite vertices of these
cycles.

\medskip\noindent\emph{Step 3: the neighbours in $X_2$ of a vertex of $X_0$.}
Let $z\in X_0$. It has no neighbour in $P\cup Q$. If $z$ were adjacent to two linked vertices of $X_2$, these two
distinct vertices would have three common neighbours: two in $P\cup Q$ and $z$. So in each of the two cycles the
neighbours of $z$ are pairwise not linked: there are $r_z\le2$ of them in the first cycle and $s_z\le2$ in the
second, and when there are two they form a diagonal. Let $q_z$ be the number of diagonals both of whose vertices are
adjacent to $z$. Then $q_z=\binom{r_z}2+\binom{s_z}2$, and since $\binom r2\ge r-1$ for $r\in\{0,1,2\}$,
\[
q_z\ \ge\ r_z+s_z-2=\alpha_z-2 .
\]

\medskip\noindent\emph{Step 4: contradiction.}
By Corollary~\ref{cor:ten}, $\sum_{z\in X_0}q_z\ge\sum_{z\in X_0}(\alpha_z-2)\ge10$. But there are four diagonals,
and the two vertices of a diagonal are distinct, so they have at most two common neighbours:
$\sum_{z\in X_0}q_z\le4\cdot2=8$.
\end{proof}

\begin{proof}[Proof of Theorem~\ref{thm:A}]
(a) is Lemma~\ref{lem:values}, (b) is Lemma~\ref{lem:indep}, (c) is Lemma~\ref{lem:k0}. For (d): by
Lemma~\ref{lem:C} and the remark before it, $\mathcal C$ is a graph on $8$ vertices, regular of degree $3$, without
triangle; by Lemma~\ref{lem:cubic} it is the cube or the Wagner graph; by Proposition~\ref{prop:cube} it is not the
cube.
\end{proof}

\end{document}
